\documentclass[10pt,a4paper]{amsart}
\usepackage{amsmath,amssymb,amsthm,graphicx}
% Retained preamble for the Stacks Project chapter morphisms.tex (Morphisms of Schemes).
% The chapter file pinned-original-source.tex carries only \input{preamble} on its first line, so no
% preamble text lives inside the chapter itself.  The Stacks Project uses one single shared preamble for
% every chapter; the snapshot stages it verbatim in sources/stacks-chow/preamble.tex.  The block below is
% copied mechanically, line for line, from that file's lines 14-18 (the Xy-pic setup) and lines 154-217
% (the native theorem declarations together with the mathematical macros).  Nothing was added, removed or
% redefined.  The wrapper class and the amsmath/amssymb/amsthm/graphicx base packages are supplied by the
% assembler, so the class/\externaldocument part of the shared preamble is not repeated here.
% For commutative diagrams we use Xy-pic
\usepackage[all]{xy}

% We use 2cell for 2-commutative diagrams.
\xyoption{2cell}
% Theorem environments.
%
\theoremstyle{plain}
\newtheorem{theorem}[subsection]{Theorem}
\newtheorem{proposition}[subsection]{Proposition}
\newtheorem{lemma}[subsection]{Lemma}

\theoremstyle{definition}
\newtheorem{definition}[subsection]{Definition}
\newtheorem{example}[subsection]{Example}
\newtheorem{exercise}[subsection]{Exercise}
\newtheorem{situation}[subsection]{Situation}

\theoremstyle{remark}
\newtheorem{remark}[subsection]{Remark}
\newtheorem{remarks}[subsection]{Remarks}

\numberwithin{equation}{subsection}

% Macros
%
\def\lim{\mathop{\mathrm{lim}}\nolimits}
\def\colim{\mathop{\mathrm{colim}}\nolimits}
\def\Spec{\mathop{\mathrm{Spec}}}
\def\Hom{\mathop{\mathrm{Hom}}\nolimits}
\def\Ext{\mathop{\mathrm{Ext}}\nolimits}
\def\SheafHom{\mathop{\mathcal{H}\!\mathit{om}}\nolimits}
\def\SheafExt{\mathop{\mathcal{E}\!\mathit{xt}}\nolimits}
\def\Sch{\mathit{Sch}}
\def\Mor{\mathop{\mathrm{Mor}}\nolimits}
\def\Ob{\mathop{\mathrm{Ob}}\nolimits}
\def\Sh{\mathop{\mathit{Sh}}\nolimits}
\def\NL{\mathop{N\!L}\nolimits}
\def\CH{\mathop{\mathrm{CH}}\nolimits}
\def\proetale{{pro\text{-}\acute{e}tale}}
\def\etale{{\acute{e}tale}}
\def\QCoh{\mathit{QCoh}}
\def\Ker{\mathop{\mathrm{Ker}}}
\def\Im{\mathop{\mathrm{Im}}}
\def\Coker{\mathop{\mathrm{Coker}}}
\def\Coim{\mathop{\mathrm{Coim}}}

% Boxtimes
%
\DeclareMathSymbol{\boxtimes}{\mathbin}{AMSa}{"02}

%
% Macros for moduli stacks/spaces
%
\def\QCohstack{\mathcal{QC}\!\mathit{oh}}
\def\Cohstack{\mathcal{C}\!\mathit{oh}}
\def\Spacesstack{\mathcal{S}\!\mathit{paces}}
\def\Quotfunctor{\mathrm{Quot}}
\def\Hilbfunctor{\mathrm{Hilb}}
\def\Curvesstack{\mathcal{C}\!\mathit{urves}}
\def\Polarizedstack{\mathcal{P}\!\mathit{olarized}}
\def\Complexesstack{\mathcal{C}\!\mathit{omplexes}}
% \Pic is the operator that assigns to X its picard group, usage \Pic(X)
% \Picardstack_{X/B} denotes the Picard stack of X over B
% \Picardfunctor_{X/B} denotes the Picard functor of X over B
\def\Pic{\mathop{\mathrm{Pic}}\nolimits}
\def\Picardstack{\mathcal{P}\!\mathit{ic}}
\def\Picardfunctor{\mathrm{Pic}}
\def\Deformationcategory{\mathcal{D}\!\mathit{ef}}

% --- appended 2026-10-12 (Chief): the project own definitions of the reference/slogan/history
% display environments, copied mechanically from the shared Stacks preamble lines 8-12
% (\usepackage{verbatim} plus the three \newenvironment lines, which the chapter preamble keeps
% verbatim from the project preamble). Nothing added, removed or redefined: these environments
% carry no mathematical content and the project itself typesets them as comments. Without these
% three lines a counted statement whose own text uses them fails with Environment slogan undefined.
\usepackage{verbatim}
\newenvironment{reference}{\comment}{\endcomment}
\newenvironment{slogan}{\comment}{\endcomment}
\newenvironment{history}{\comment}{\endcomment}
\title{An A-module on a scheme with quasi-coherent algebra A is quasi-coherent over A exactly when it is quasi-coherent over the structure sheaf}
\author{The Stacks Project Authors}
\date{stacks-project:morphisms.tex:local-snapshot-20261012; GNU FDL 1.2 or later}
\begin{document}
\maketitle
\noindent\textit{Editorial note.} All delivered mathematical text is the author's own, copied from the pinned Stacks Project chapter file pinned-original-source.tex; no statement or proof step was rewritten, reordered or removed.\par\medskip
\noindent\textit{Editorial note.} Following the multi-claim convention for this source, source\_url is the source-level URL https://stacks.math.columbia.edu/ and the source label lemma-quasi-coherence-scalar-restriction is carried in claim\_key and here rather than in a per-claim tag URL.\par\medskip
\noindent\textit{Editorial note.} The preamble is the chapter's shared Stacks Project preamble, retained as sources/stacks-morphisms/retained-preamble.tex: a line-for-line excerpt of the shared Stacks preamble staged with this snapshot. No macro was invented or redefined, and this card adds no preamble text of its own.\par\medskip
\noindent\textit{Editorial note.} Statement span 1844-1853 and proof span 1855-1893 come from the halfA pool card operations/pool-cards/stacks-morphisms-halfA.json and were re-read in the pinned chapter file before this card was written.\par\medskip
\noindent\textit{Editorial note.} The counted statement contains a reference environment holding the author's citation of EGA. The project's own shared preamble (now retained in full, 88 lines, per the 2026-10-12 repair) defines the reference environment as a comment environment, so that citation is not typeset by the project itself either; a bibliography is therefore not attached to this card and none is needed.\par\medskip
\noindent\textit{Editorial note.} The counted proof contains the author's parenthetical '(details omitted)' for the routine point that the module structure on M is compatible with the R-module structure; that phrase is kept verbatim and what it leaves to the reader is a compatibility check, not a step of the main argument.\par\medskip
\noindent\textit{Allowed standard prerequisites (editorial).}\begin{itemize}
\item Local presentations of a quasi-coherent module as the cokernel of a map of direct sums of the structure sheaf or of the algebra, on an open neighbourhood of a point. This is what the two directions of the proof use at proof lines 2074-2078 and 2095-2102. Supplies the exact sequences in terms of which the module is recognised as quasi-coherent over the algebra, respectively over the structure sheaf.
\item The tilde construction over an affine scheme: a quasi-coherent algebra on Spec(R) is the tilde of an R-algebra, a quasi-coherent module is the tilde of an R-module, and the tilde functor is exact. Used at proof lines 2086-2092 and again at 2103-2108. Translates the second direction into a statement about an R-algebra A and an R-module M with an A-module structure, and transfers an A-module presentation back to sheaves.
\item Cokernels of maps of quasi-coherent modules on a scheme are quasi-coherent, and exactness of a sequence of modules over an algebra regarded as a sequence of modules over the smaller ring. Used at proof lines 2077-2080 and 2092-2093. Closes the first direction and identifies the compatible structures in the affine chart.
\end{itemize}\medskip
\begin{lemma}
\label{lemma-quasi-coherence-scalar-restriction}
\begin{reference}
\cite[I, Proposition 9.6.1]{EGA}
\end{reference}
Let $S$ be a scheme and let $\mathcal{A}$ be a quasi-coherent
$\mathcal{O}_S$-algebra. An $\mathcal{A}$-module is
quasi-coherent as an $\mathcal{O}_S$-module if and only if
it is quasi-coherent as an $\mathcal{A}$-module.
\end{lemma}

\begin{proof}
Let $\mathcal{F}$ be an $\mathcal{A}$-module.
If $\mathcal{F}$ is quasi-coherent as an $\mathcal{A}$-module, then
for every $s \in S$ there exists an open neighbourhood $U$ of $s$
and an exact sequence
$$
\bigoplus\nolimits_J\mathcal{A}|_U\to
\bigoplus\nolimits_I\mathcal{A}|_U\to
\mathcal{F}|_U
$$
of $\mathcal{A}|_U$-modules.
Then this is also an exact sequence of $\mathcal{O}_U$-modules.
Hence $\mathcal{F}|_U$ is quasi-coherent as the cokernel of a morphism
of quasi-coherent $\mathcal{O}_U$-modules on a scheme. It follows
that $\mathcal{F}$ is quasi-coherent as an $\mathcal{O}_X$-module.

\medskip\noindent
Conversely, assume $\mathcal{F}$ is quasi-coherent as an $\mathcal{O}_X$-module.
Pick an open affine $\Spec(R) = U \subset S$. We have isomorphisms of
$\mathcal{O}_U$-modules $\mathcal{A}|_U \cong \widetilde{A}$ and
$\mathcal{F}|_U \cong \widetilde{M}$, for some $R$-algebra $A$
and some $R$-module $M$. The $\mathcal{A}$-module structure on $\mathcal{F}$
translates into an $A$-module structure on $M$ compatible with the
given $R$-module structure (details omitted). Choose an exact sequence
$$
\bigoplus\nolimits_J A \to
\bigoplus\nolimits_I A \to
M \to 0
$$
of $A$-modules. Since the functor $\widetilde{\ }$ is exact, this
produces an exact sequence
$$
\bigoplus\nolimits_J \widetilde{A}\to
\bigoplus\nolimits_I \widetilde{A}\to
\widetilde{M} \to 0
$$
of $\widetilde{A}$-modules. This means that
$\mathcal{F}$ is quasi-coherent as an $\mathcal{A}$-module.
\end{proof}

\end{document}
